% Options for packages loaded elsewhere
\PassOptionsToPackage{unicode}{hyperref}
\PassOptionsToPackage{hyphens}{url}
\PassOptionsToPackage{dvipsnames,svgnames*,x11names*}{xcolor}
%
\documentclass[
  11pt,
]{article}
\usepackage{lmodern}
\usepackage{amssymb,amsmath}
\usepackage{ifxetex,ifluatex}
\ifnum 0\ifxetex 1\fi\ifluatex 1\fi=0 % if pdftex
  \usepackage[T1]{fontenc}
  \usepackage[utf8]{inputenc}
  \usepackage{textcomp} % provide euro and other symbols
\else % if luatex or xetex
  \usepackage{unicode-math}
  \defaultfontfeatures{Scale=MatchLowercase}
  \defaultfontfeatures[\rmfamily]{Ligatures=TeX,Scale=1}
\fi
% Use upquote if available, for straight quotes in verbatim environments
\IfFileExists{upquote.sty}{\usepackage{upquote}}{}
\IfFileExists{microtype.sty}{% use microtype if available
  \usepackage[]{microtype}
  \UseMicrotypeSet[protrusion]{basicmath} % disable protrusion for tt fonts
}{}
\makeatletter
\@ifundefined{KOMAClassName}{% if non-KOMA class
  \IfFileExists{parskip.sty}{%
    \usepackage{parskip}
  }{% else
    \setlength{\parindent}{0pt}
    \setlength{\parskip}{6pt plus 2pt minus 1pt}}
}{% if KOMA class
  \KOMAoptions{parskip=half}}
\makeatother
\usepackage{xcolor}
\IfFileExists{xurl.sty}{\usepackage{xurl}}{} % add URL line breaks if available
\IfFileExists{bookmark.sty}{\usepackage{bookmark}}{\usepackage{hyperref}}
\hypersetup{
  colorlinks=true,
  linkcolor=blue,
  filecolor=Maroon,
  citecolor=Blue,
  urlcolor=blue,
  pdfcreator={LaTeX via pandoc}}
\urlstyle{same} % disable monospaced font for URLs
\usepackage[margin=1in]{geometry}
\usepackage{longtable,booktabs}
% Correct order of tables after \paragraph or \subparagraph
\usepackage{etoolbox}
\makeatletter
\patchcmd\longtable{\par}{\if@noskipsec\mbox{}\fi\par}{}{}
\makeatother
% Allow footnotes in longtable head/foot
\IfFileExists{footnotehyper.sty}{\usepackage{footnotehyper}}{\usepackage{footnote}}
\makesavenoteenv{longtable}
\setlength{\emergencystretch}{3em} % prevent overfull lines
\providecommand{\tightlist}{%
  \setlength{\itemsep}{0pt}\setlength{\parskip}{0pt}}
\setcounter{secnumdepth}{-\maxdimen} % remove section numbering

\author{}
\date{}

\begin{document}

\hypertarget{calibration-free-fp8-quantization-for-transformers-a-shannon-prime-experiment-design}{%
\section{Calibration-Free fp8 Quantization for Transformers: A
Shannon-Prime Experiment
Design}\label{calibration-free-fp8-quantization-for-transformers-a-shannon-prime-experiment-design}}

\textbf{A. Knack} (Shannon-Prime Project) \textbf{Companion systems
paper. Draft v0.1 --- 2026-05-16}

\begin{center}\rule{0.5\linewidth}{0.5pt}\end{center}

\hypertarget{abstract}{%
\subsection{Abstract}\label{abstract}}

This companion paper specifies the experiment that probes the single
highest-leverage prediction of the Shannon-Prime framework: that fp8
quantization on a CM-encoded transformer state requires no calibration.
The prediction is a direct consequence of the Frobenius Quantization
Theorem (Theorem 2 of the theory companion {[}SP-Theory-Companion
2026{]}); this paper presents the experimental design that tests it on a
Phi-3 reference model, with the full implementation hooks already
present in the Shannon-Prime engine. We describe four configurations
(baseline fp16, SP-fp8 calibration-free, GPTQ-fp8 calibrated, AWQ-fp8
calibrated) and the evaluation harness, perplexity targets, and
pass/fail criteria. The expected outcome is that SP-fp8 matches the
calibrated baselines without any calibration data. We include
implementation notes, a reproducible run script outline, and a
discussion of how the result, if positive, scales to fp4.

\begin{center}\rule{0.5\linewidth}{0.5pt}\end{center}

\hypertarget{the-question}{%
\subsection{1. The Question}\label{the-question}}

Standard transformer quantization treats the model's weight and
activation tensors as floating-point arrays and rounds them to a
lower-precision representation. The rounding introduces errors in the
multiplicative relations the model has learned during training, and
post-training calibration is required to fit a per-tensor scale factor
that approximately restores those relations. GPTQ, AWQ, SmoothQuant, and
QServe all follow this pattern; each requires a calibration pass on
representative data (typically \textasciitilde1000 samples) to recover
acceptable accuracy.

Shannon-Prime predicts something different. Its theoretical content
(Theorem 2 of {[}SP-Theory-Companion 2026{]}) is that on a CM elliptic
curve over \(\mathbb{Q}(\sqrt{-163})\), the Frobenius endomorphism
\(\varphi_p\) commutes with the layer endomorphisms. If the model's
state is encoded on this curve --- as SP encodes the KV cache and (in
this experiment) the weights --- then reduction to lower precision is
Frobenius application rather than rounding, and the multiplicative
relations are preserved exactly.

\textbf{Question.} \emph{Does SP-fp8 on a Phi-3 model recover
vanilla-fp16 perplexity to within noise, with no calibration data?}

A positive answer would deliver one of the strongest practical results
in the field: a quantization scheme requiring no calibration pass, hence
shippable to edge devices without device-specific data.

\begin{center}\rule{0.5\linewidth}{0.5pt}\end{center}

\hypertarget{setup}{%
\subsection{2. Setup}\label{setup}}

\hypertarget{reference-model}{%
\subsubsection{2.1 Reference Model}\label{reference-model}}

Phi-3 was chosen for three reasons:

\begin{enumerate}
\def\labelenumi{\arabic{enumi}.}
\tightlist
\item
  \textbf{It is the SP model-pack's calibrated reference.} Phi-3
  currently passes the SP calibration ledger at
  \(\Delta\mathrm{PPL} = +2.44\%\) vs vanilla under SP-base compression.
  Starting from a known-calibrated SP target removes one source of
  variance.
\item
  \textbf{It is small enough to bench quickly.} 3.8B parameters at fp16
  fits on a single A100; the bench runs in under an hour per
  configuration.
\item
  \textbf{It does not exhibit the Qwen3 edge-fail artifact.} The Qwen3
  architecture's gated-attention + mRoPE-mode-8 combination breaks the
  SP encoding in a known way; Phi-3 does not. Using Phi-3 eliminates
  that confound.
\end{enumerate}

\hypertarget{bench-corpora}{%
\subsubsection{2.2 Bench Corpora}\label{bench-corpora}}

Two corpora at three context lengths each:

\begin{longtable}[]{@{}lll@{}}
\toprule
Corpus & Description & Tokens\tabularnewline
\midrule
\endhead
\textbf{WikiText-103} & Standard perplexity bench & 245M\tabularnewline
\textbf{The Stack v2 (filtered)} & Code completion bench &
100M\tabularnewline
\bottomrule
\end{longtable}

Context lengths: 512, 2048, 8192. The 8192 condition exercises the
long-context KV cache path where SP compression is most impactful.

\hypertarget{configurations}{%
\subsubsection{2.3 Configurations}\label{configurations}}

Four configurations, all evaluated on the same corpus + context
combinations:

\begin{longtable}[]{@{}llll@{}}
\toprule
Config & Quant & Calibration & Notes\tabularnewline
\midrule
\endhead
\textbf{A. Baseline} & fp16 & --- & Ground truth\tabularnewline
\textbf{B. SP-fp8 calibration-free} & SP-fp8 (\(p = 11\),
\(\varphi_p^8\)) & None & The hypothesis\tabularnewline
\textbf{C. GPTQ-fp8} & fp8 & \textasciitilde1000 samples & Industry
baseline\tabularnewline
\textbf{D. AWQ-fp8} & fp8 & \textasciitilde1000 samples & Industry
baseline\tabularnewline
\bottomrule
\end{longtable}

The configuration that is the focus is B. The hypothesis is that B
matches A within noise (\(|\Delta\mathrm{PPL}| < 0.5\%\)) and matches
C/D within noise (\(|\Delta\mathrm{PPL}_{B,C}| < 0.3\%\)).

\begin{center}\rule{0.5\linewidth}{0.5pt}\end{center}

\hypertarget{why-fp8-why-p-11-why-varphi_p8}{%
\subsection{\texorpdfstring{3. Why fp8, Why \(p = 11\), Why
\(\varphi_p^8\)}{3. Why fp8, Why p = 11, Why \textbackslash varphi\_p\^{}8}}\label{why-fp8-why-p-11-why-varphi_p8}}

The Frobenius Quantization Theorem (Theorem 2 of {[}SP-Theory-Companion
2026{]}) implements \(\mathrm{fp}16 \to \mathrm{fp}q\) as iterated
Frobenius \(\varphi_p^{16-q}\). For \(q = 8\), this is \(\varphi_p^8\).
The prime \(p\) is constrained by \(p^q \leq 2^{16}\), i.e.,
\(p^8 \leq 2^{16}\), giving \(p \leq 2^2 = 4\), which is too small to be
useful at \(q = 8\).

The resolution is to relax the constraint: at \(q = 8\) we have an
entire byte of representation, and \(p\) can be chosen up to
\(2^8 = 256\). Setting \(p = 11\) gives a comfortable margin while
keeping the trace of Frobenius \(a_{11}\) in a numerically friendly
range. For \(E\) the CM curve over \(\mathbb{Q}(\sqrt{-163})\) reduced
mod 11, \(a_{11}\) can be computed from the Hilbert class polynomial;
the explicit value is \(a_{11} = -6\) (good ordinary reduction). The
Frobenius polynomial is then \(\varphi_{11}^2 + 6\varphi_{11} + 11 = 0\)
in \(\operatorname{End}(E_{11})\).

The number of points:
\(\#E_{11}(\mathbb{F}_{11}) = 11 + 1 - a_{11} = 18\), well within the
Hasse--Weil bound of \(11 + 1 + 2\sqrt{11} \approx 18.63\).

\textbf{At \(q = 8\), per-coordinate information capacity is
\(\log_2 18 \approx 4.17\) bits.} This is lower than the nominal ``8
bits'' of fp8 because the CM encoding uses the upper bits for spinor and
Möbius coefficients. Empirically this matches the effective bit-width of
well-quantized fp8 LLM weights as reported in {[}Liu et al.~2024
KIVI{]}.

\begin{center}\rule{0.5\linewidth}{0.5pt}\end{center}

\hypertarget{implementation}{%
\subsection{4. Implementation}\label{implementation}}

\hypertarget{existing-hooks}{%
\subsubsection{4.1 Existing Hooks}\label{existing-hooks}}

The Shannon-Prime engine already contains:

\begin{itemize}
\tightlist
\item
  \texttt{-\/-hier-ternary-mask} and \texttt{-\/-hier-res-bits-v} flags
  for ternary skeleton compression.
\item
  VHT2 + Möbius + spinor + square-free chain across all four backends.
\item
  The fp8 quantization implementation in the engine path (used in
  production but currently calibrated).
\item
  Per-architecture compression-default registry (model-pack scaffold)
  --- phi3 row already exists and is calibrated.
\end{itemize}

What is needed to run the experiment:

\begin{enumerate}
\def\labelenumi{\arabic{enumi}.}
\tightlist
\item
  \textbf{A \texttt{-\/-frobenius-quant} flag} that toggles the
  Frobenius reduction in place of the calibrated fp8 path.
\item
  \textbf{A unit test} verifying that \(\varphi_{11}^8\) produces output
  bit-identical to direct reduction mod \(11^8\) on the CM-encoded
  state.
\item
  \textbf{A bench harness modification} to disable the calibration-data
  path for config B.
\end{enumerate}

Estimated implementation: 1--2 days at the engine level. The math is
fully specified by Theorem 2.

\hypertarget{outline-of-the-run-script}{%
\subsubsection{4.2 Outline of the Run
Script}\label{outline-of-the-run-script}}

\begin{verbatim}
# Pseudocode for the bench harness
for cfg in [A_baseline_fp16, B_sp_fp8_calfree, C_gptq_fp8, D_awq_fp8]:
    model = load_phi3(cfg.quantization)
    for corpus in [wikitext103, stackv2]:
        for ctx in [512, 2048, 8192]:
            ppl = eval_perplexity(model, corpus, ctx)
            tokens_per_sec = bench_throughput(model, corpus, ctx)
            log(cfg, corpus, ctx, ppl, tokens_per_sec)
\end{verbatim}

All four configs use the same bench loader and the same tokenization;
only the quantization differs.

\hypertarget{passfail-criteria}{%
\subsubsection{4.3 Pass/Fail Criteria}\label{passfail-criteria}}

\begin{longtable}[]{@{}ll@{}}
\toprule
\begin{minipage}[b]{0.47\columnwidth}\raggedright
Comparison\strut
\end{minipage} & \begin{minipage}[b]{0.47\columnwidth}\raggedright
Criterion\strut
\end{minipage}\tabularnewline
\midrule
\endhead
\begin{minipage}[t]{0.47\columnwidth}\raggedright
\textbf{B vs A} (calibration-free fp8 vs baseline fp16)\strut
\end{minipage} & \begin{minipage}[t]{0.47\columnwidth}\raggedright
\(|\Delta\mathrm{PPL}| < 0.5\%\) on WikiText-103 at all contexts\strut
\end{minipage}\tabularnewline
\begin{minipage}[t]{0.47\columnwidth}\raggedright
\textbf{B vs C} (SP-fp8 vs GPTQ-fp8)\strut
\end{minipage} & \begin{minipage}[t]{0.47\columnwidth}\raggedright
\(|\Delta\mathrm{PPL}_{B,C}| < 0.3\%\) at all contexts\strut
\end{minipage}\tabularnewline
\begin{minipage}[t]{0.47\columnwidth}\raggedright
\textbf{B vs D} (SP-fp8 vs AWQ-fp8)\strut
\end{minipage} & \begin{minipage}[t]{0.47\columnwidth}\raggedright
\(|\Delta\mathrm{PPL}_{B,D}| < 0.3\%\) at all contexts\strut
\end{minipage}\tabularnewline
\begin{minipage}[t]{0.47\columnwidth}\raggedright
\textbf{B throughput}\strut
\end{minipage} & \begin{minipage}[t]{0.47\columnwidth}\raggedright
within 5\% of A throughput; SP-fp8 is not expected to be faster than
calibrated fp8\strut
\end{minipage}\tabularnewline
\begin{minipage}[t]{0.47\columnwidth}\raggedright
\textbf{Robustness}\strut
\end{minipage} & \begin{minipage}[t]{0.47\columnwidth}\raggedright
B's relative PPL gap does not increase with context length\strut
\end{minipage}\tabularnewline
\bottomrule
\end{longtable}

If all four criteria pass, the framework's central practical prediction
is validated and SP-fp8 is deployable as a calibration-free
quantization.

\begin{center}\rule{0.5\linewidth}{0.5pt}\end{center}

\hypertarget{expected-outcome-and-risk-analysis}{%
\subsection{5. Expected Outcome and Risk
Analysis}\label{expected-outcome-and-risk-analysis}}

\hypertarget{expected-outcome}{%
\subsubsection{5.1 Expected Outcome}\label{expected-outcome}}

Theorem 2 predicts that B matches A within numerical noise. The
framework attributes any residual gap to:

\begin{itemize}
\tightlist
\item
  The Frobenius trace \(a_p\) being nonzero (predicted systematic drift
  of \(\approx a_p / p \approx 0.5\%\) at \(p = 11\)); the experiment
  can quantify this.
\item
  Bench-scaffold artifacts (the fp32 reduction step in the engine bench
  loop introduces small additional error; see the A100 result in
  {[}SP-Systems 2026{]}).
\item
  Tokenizer interaction with the CRT-sharded LM head (not exercised by
  config A but exercised by SP encoding).
\end{itemize}

The expected magnitude of \(|\Delta\mathrm{PPL}_{B,A}|\) is in the range
\(0.1\%\) to \(0.5\%\). Anything below \(1\%\) is a positive result;
anything below \(0.3\%\) is a strong positive.

\hypertarget{what-a-negative-result-would-indicate}{%
\subsubsection{5.2 What a Negative Result Would
Indicate}\label{what-a-negative-result-would-indicate}}

A failure of B to match A would localize one of the following:

\begin{enumerate}
\def\labelenumi{\arabic{enumi}.}
\tightlist
\item
  The CM realization of Phi-3's layer endomorphisms is imperfect --- the
  Endomorphism Realization (Theorem 1 of {[}SP-Theory 2026{]}) holds in
  principle but does not yet match the empirical model in
  implementation.
\item
  The Frobenius reduction is implemented incorrectly in the engine.
\item
  The framework's choice of curve \(E\) over \(\mathbb{Q}(\sqrt{-163})\)
  does not match the actual hidden-state geometry of trained Phi-3.
\end{enumerate}

\begin{enumerate}
\def\labelenumi{(\arabic{enumi})}
\tightlist
\item
  is debuggable by mechanistic interpretability of layer endomorphisms;
  (2) is a code review; (3) would require revisiting the framework's
  curve choice. Of these, (3) is the lowest-probability and the most
  consequential.
\end{enumerate}

\hypertarget{what-a-positive-result-enables}{%
\subsubsection{5.3 What a Positive Result
Enables}\label{what-a-positive-result-enables}}

A positive result delivers:

\begin{enumerate}
\def\labelenumi{\arabic{enumi}.}
\tightlist
\item
  \textbf{Calibration-free fp8 deployment.} Drop-in fp8 for any
  Phi-3-class model with no calibration corpus required.
\item
  \textbf{A bridge to fp4.} Theorem 2 with \(q = 4\) requires
  \(p \leq 11\), and the SP-fp4 numerical information ceiling of
  \(\approx 4.07\) bits/coord (Corollary 2.2 of the theory companion)
  matches the empirical fp4 capacity. The same experiment, repeated at
  \(q = 4\) with appropriate \(p\), would be the natural next step.
\item
  \textbf{A target for ARM/Hexagon production.} Phone-class targets
  currently quantize via runtime-calibrated routines. Calibration-free
  quantization removes a major deployment friction.
\end{enumerate}

\begin{center}\rule{0.5\linewidth}{0.5pt}\end{center}

\hypertarget{implementation-timeline}{%
\subsection{6. Implementation Timeline}\label{implementation-timeline}}

\begin{longtable}[]{@{}lll@{}}
\toprule
\begin{minipage}[b]{0.30\columnwidth}\raggedright
Phase\strut
\end{minipage} & \begin{minipage}[b]{0.30\columnwidth}\raggedright
Deliverable\strut
\end{minipage} & \begin{minipage}[b]{0.30\columnwidth}\raggedright
Estimated effort\strut
\end{minipage}\tabularnewline
\midrule
\endhead
\begin{minipage}[t]{0.30\columnwidth}\raggedright
1. Engine hooks\strut
\end{minipage} & \begin{minipage}[t]{0.30\columnwidth}\raggedright
\texttt{-\/-frobenius-quant} flag, unit test for
\(\varphi_{11}^8\)\strut
\end{minipage} & \begin{minipage}[t]{0.30\columnwidth}\raggedright
1--2 days\strut
\end{minipage}\tabularnewline
\begin{minipage}[t]{0.30\columnwidth}\raggedright
2. Bench harness\strut
\end{minipage} & \begin{minipage}[t]{0.30\columnwidth}\raggedright
Modified eval loop for the four configs\strut
\end{minipage} & \begin{minipage}[t]{0.30\columnwidth}\raggedright
1 day\strut
\end{minipage}\tabularnewline
\begin{minipage}[t]{0.30\columnwidth}\raggedright
3. Reference runs\strut
\end{minipage} & \begin{minipage}[t]{0.30\columnwidth}\raggedright
All four configs at all six (corpus × context) combinations\strut
\end{minipage} & \begin{minipage}[t]{0.30\columnwidth}\raggedright
1 day on A100\strut
\end{minipage}\tabularnewline
\begin{minipage}[t]{0.30\columnwidth}\raggedright
4. Analysis + writeup\strut
\end{minipage} & \begin{minipage}[t]{0.30\columnwidth}\raggedright
Result table, plot of PPL vs context, regression to the framework's
predictions\strut
\end{minipage} & \begin{minipage}[t]{0.30\columnwidth}\raggedright
2 days\strut
\end{minipage}\tabularnewline
\begin{minipage}[t]{0.30\columnwidth}\raggedright
5. Companion v0.2\strut
\end{minipage} & \begin{minipage}[t]{0.30\columnwidth}\raggedright
Incorporate results into this paper as §5 (Results)\strut
\end{minipage} & \begin{minipage}[t]{0.30\columnwidth}\raggedright
1 day\strut
\end{minipage}\tabularnewline
\bottomrule
\end{longtable}

Total: \textasciitilde1 week of engine + bench work, assuming an A100 is
available for the reference runs.

\begin{center}\rule{0.5\linewidth}{0.5pt}\end{center}

\hypertarget{connection-to-the-broader-sp-framework}{%
\subsection{7. Connection to the Broader SP
Framework}\label{connection-to-the-broader-sp-framework}}

This experiment is one component of a larger program. The Shannon-Prime
framework predicts five additional implementations, each with its own
experimental signature:

\begin{itemize}
\tightlist
\item
  \textbf{Stern--Brocot RoPE} ({[}SP-Theory 2026, §9.1{]}): replace RoPE
  base \(10{,}000\) with \(\varphi\); bench on long-context tasks.
  Smallest implementation cost of any extension.
\item
  \textbf{Weil pairing attention} ({[}SP-Theory 2026, §9.2{]}): replace
  softmax-of-dot-product with the Weil pairing \(e_n\). Highest research
  upside.
\item
  \textbf{Hecke-eigenform embeddings} ({[}SP-Theory 2026, §9.3{]}):
  replace learned embedding with a Hecke basis. Longest path to a
  number.
\item
  \textbf{L-function activation oracle} ({[}SP-Theory 2026, §9.4{]}):
  predict FFN firing as L-function coefficients.
\item
  \textbf{LLL KV write} ({[}SP-Theory 2026, §9.5{]}): make the KV
  archive an LLL-reduced lattice basis.
\end{itemize}

Each of these is the subject of a future SP companion paper. The present
experiment --- calibration-free fp8 --- is first in line because it is
the smallest implementation step that probes the strongest theoretical
prediction.

\begin{center}\rule{0.5\linewidth}{0.5pt}\end{center}

\hypertarget{conclusion}{%
\subsection{8. Conclusion}\label{conclusion}}

Shannon-Prime's Frobenius Quantization Theorem predicts that fp8
quantization on a CM-encoded transformer state requires no calibration
data. The present companion paper has specified the experimental design
that tests this prediction on a Phi-3 reference model. The result, if
positive, validates the framework's strongest practical claim and opens
a path to fp4 and edge-device deployment without calibration. The
experiment is small (≈1 week of engineering + benchmarking), the
implementation hooks are already in place across the SP engine and
llama.cpp fork, and the framework's predicted outcome is sharp and
falsifiable.

\begin{center}\rule{0.5\linewidth}{0.5pt}\end{center}

\hypertarget{references}{%
\subsection{References}\label{references}}

\begin{itemize}
\tightlist
\item
  Companion theory: ``Two Theorems for Shannon-Prime Compression:
  Hasse--Weil as the Shannon Limit, and Frobenius as Quantization.''
  2026.
\item
  Full theory: ``Shannon-Prime: A CM-Elliptic-Curve Framework for
  Transformer Computation.'' 2026.
\item
  Full systems: ``Shannon-Prime: Provably Exact KV-Cache Compression and
  Per-Layer Acceleration on Standard Inference Stacks.'' 2026.
\item
  Liu, Z. \emph{et al.} ``KIVI: A Tuning-Free Asymmetric 2bit
  Quantization for KV Cache.'' 2024.
\item
  Frantar, E. \emph{et al.} ``GPTQ: Accurate Post-Training Quantization
  for Generative Pre-trained Transformers.'' 2023.
\item
  Lin, J. \emph{et al.} ``AWQ: Activation-aware Weight Quantization for
  LLM Compression and Acceleration.'' 2024.
\item
  Phi-3 Technical Report, Microsoft, 2024.
\end{itemize}

\end{document}
